\section{Background}

Large Language Models (LLMs) have displayed exceptional proficiency across various NLP (Natural Language Processing) tasks such as answering questions, summarizing, generating codes and reasoning. Models such as GPT-4, Claude and PaLM achieve impressive performance on usual benchmarks, with in some cases, even exceeding human-level accuracy on certain tasks \cite{openai2023gpt4, anthropic2023claude, anil2023palm}. However, a critical limitation hampers their deployment in high-stakes applications, for example, the generation of plausible yet factually incorrect responses, commonly known as \textit{hallucinations}.

In LLMs, hallucination occurs when models produce outputs that contradict source material, conflict with valid world knowledge, or even fabricate information entirely. Recent studies reveal concerning statistics about hallucination occurrences. The TruthfulQA benchmark, introduced by Lin et al. \cite{lin-etal-2022-truthfulqa}, is designed to particularly measure how models copy human falsehoods. The worst part is that it found larger models performing worse than small ones. For example, GPT-3 175B achieved only 58\% on truthfulness alone and just 21--25\% on the combined truthful-and-informative metric, compared to 94\% human accuracy on both metrics, demonstrating an \textit{inverse scaling} phenomenon where increased model size does not guarantee improved truthfulness. Also, another benchmark, FActScore by Min et al. \cite{min-etal-2023-factscore}, developed for fine-grained evaluation of long-form generation, shows that ChatGPT achieved only 58\% supported atomic facts in biography generation. Domain-specific studies show even higher error rates. For example, medical literature analysis by Alkaissi and McFarlane \cite{alkaissi2023artificial} found hallucination rates of 69-77\% in ChatGPT responses to clinical questions.

The fundamental challenge expands beyond accuracy benchmarks. Current LLMs operate statelessly where each query is processed independently without learning from previous interactions or collecting verified knowledge. This stateless architecture results in three critical limitations. Firstly, \textit{no memory of successful reasoning patterns}, which requires generation from scratch for every query. Secondly, \textit{poor confidence calibration}, where models show high confidence even while being incorrect \cite{kadavath2022language}. Lastly, \textit{no self-improvement mechanism}, which causes indefinite repetition of identical errors without adapting.

All these limitations have significant practical consequences. For example, in educational sectors, students may learn incorrect information presented confidently. In healthcare sectors, hallucinated medical information can lead to harmful decisions. In enterprise knowledge management, fake citations and sources can hamper trust. The stateless nature also causes the repetition of inefficient-expensive generation for similar queries rather than using past verified reasonings.

This research addresses these fundamental limitations through CAEM (Confidence-Aware Episodic Memory with Self-Improvement), a system that constantly reduces hallucination occurrences through three core changes. Firstly, \textit{verified episodic memory} that stores only high-quality reasoning chains which is verified through multi-layer verification. Secondly, \textit{multi-signal confidence estimation} that enables adaptive routing between retrieval and generation. Thirdly, \textit{iterative fine-tuning} which creates a closed feedback loop where verified examples improve model generation quality, which in turn produces better training data. Unlike post-hoc detection methods, which find the occurrence of hallucinations through generation, CAEM avoids hallucinations by combining external memory, confidence-gated routing, and progressive learning to reduce hallucinations by 20-30\%, while maintaining general language capabilities.

\section{Rationale of the Study or Motivation}

The motivation for this research is that there is a wide gap between what LLM are capable of and what is required when they are used. While current models are capable of writing fluent, relevant text with excellent capacity, they usually cannot distinguish between \textit{real facts} and \textit{confabulation}, which makes it difficult to apply them in domains requiring precision.

The existing approaches to decrease the number of hallucinations could be classified into three categories, with all of them possessing considerable limitations. Firstly, \textit{Retrieval-augmented generation} (RAG) methods \cite{NEURIPS2020_6b493230}, which looks up documents and uses them for generation. But, the model's answer is as good as the documents it retrieves and only depends on retrieval quality, not experience. Secondly, \textit{Prompting techniques} such as chain-of-thought \cite{NEURIPS2022_9d560961} and self-consistency \cite{wang2023selfconsistency} makes the model think more carefully and checks multiple reasoning paths, but cannot remember information later and each session starts with no knowledge stored. Thirdly, \textit{Post-hoc verification methods} \cite{manakul-etal-2023-selfcheckgpt, farquhar2024semantic} checks the answer after it is generated and can detect hallucinations after it happens, but does not teach the model what to do next time to reduce the mistakes, meaning the same hallucinations can happen again in the future.

The key insight driving this research is that hallucination reduction requires \textit{architectural intervention} rather than incremental improvements to existing methods. Specifically, three capabilities are essential but absent from current systems. First, models need \textit{episodic memory} to store and retrieve verified reasoning chains, enabling instant access to known-correct solutions rather than regenerating potentially incorrect responses. Second, models require \textit{calibrated confidence estimation} to determine when retrieval is appropriate versus when external verification is needed, preventing overconfident errors while avoiding unnecessary computation. Third, systems must support \textit{iterative self-improvement}, where verified examples create training data that improves generation quality, establishing a virtuous cycle of progressive accuracy gains.

This research is significant for both theoretical and practical reasons. Theoretically, it addresses the open question of whether imperfect verification can enable monotonic improvement in self-training systems. Prior work in semi-supervised learning demonstrates that noisy labels can degrade model performance \cite{natarajan2013learning}, yet our preliminary analysis suggests that when base model accuracy exceeds random chance, verified data purity can exceed verification accuracy itself through favorable base rate dynamics. This counterintuitive property, if validated empirically, provides theoretical justification for self-improvement under realistic verification constraints.

Practically, reducing hallucinations by even 20-30\% would make language models much more useful without hampering their overall abilities. This would allow educational systems to safely learn and remember specific institutional knowledge. Also, it would help businesses gradually build reliable domain expertise through verified interactions. In healthcare and law domains, it would grasp language understanding while maintaining strict accuracy standards. Additionally, using memory to reuse past answers would significantly lower computation costs, making AI systems more affordable and sustainable.

The undergraduate thesis scope is appropriate because the main ideas can be tested using well-known models like \textit{Flan-T5 Large}, public benchmarks like \textit{HotpotQA, StrategyQA, FEVER, TruthfulQA} and manageable computational resources like \textit{approximately 30--45 GPU hours} (8--12 hours per cycle for inference and verification across three cycles, plus 20--30 minutes per cycle for fine-tuning Flan-T5-Large, well within a standard university cluster allocation of 50 GPU hours). A 12 month timeline is enough to properly evaluate the system while focusing only on the factual and knowledge based tasks. The correctness can be clearly checked, while excluding creative tasks that are harder to evaluate objectively.

\section{Problem Statement}

Current large language models suffer from three fundamental limitations that prevent their reliable deployment in knowledge-intensive applications which requires factual accuracy:

\begin{itemize}
    \item \textbf{Stateless Operation Without Knowledge Accumulation:} LLMs process each query independently without retaining knowledge from previous interactions. Even if the model thinks carefully and gets the answer right, all the thinking is discarded once the answer is generated. So, if a similar question is asked again, the model must redo all the thinking again from the beginning. This stateless architecture causes two problems. Firstly, it wastes time as well as computing power by repeating the work. Secondly, the model might give a completely different answer next time, even for the same question. The absence of episodic memory prevents it from learning any domain-specific knowledge as each interaction is separated and does not help the model to get better over time.
    
    \item \textbf{Poor Confidence Calibration and Uncertainty Quantification:} Modern neural networks are often more confident than they should be, meaning they sound sure even when they are wrong. Guo et al. \cite{pmlr-v70-guo17a} shows that even when the model is correct only 70\% of the time, it may act like it is 90\% sure. This mismatch is dangerous because people may trust wrong answers too much. This problem shows up in two main ways. Firstly, the model gives a wrong answer very confidently, so users do not realize they should double-check it. Secondly, the model sometimes gives a correct answer but may sound unsure, which causes additional checking that was not required. If the model can not judge its own confidence properly, it can not decide if a quick lookup is enough or when deeper reasoning is needed. Current single-signal confidence methods achieve only 0.50-0.60 AUROC for telling whether an answer is right or wrong \cite{xiong2023can}, barely above random chance.
    
    \item \textbf{Absence of Self-Improvement Mechanisms:} Deployed LLMs remain static after training, unable to learn from mistakes or adapt to domain-specific patterns. If a model hallucinates, it will keep making the same mistake forever unless humans retrain it again. Trying to teach the model new things through naive fine-tuning causes a big problem called \textit{catastrophic forgetting} \cite{doi:10.1073/pnas.1611835114}, where learning new information causes 30-50\% degradation of already acquired capabilities. The absence of self-improvement means that verification efforts provide no long-term value, for example, verifying 1000 answers today provides no benefit tomorrow, as the system cannot combine verified knowledge into improved generation.
\end{itemize}

These three limitations make each other worse, creating systems that remain simultaneously overconfident and unreliable, expensive yet static, capable yet ultimately unsuitable for high-stakes deployment. This research addresses all three simultaneously through architectural innovation such as verified episodic memory providing knowledge accumulation, multi-signal confidence estimation enabling reliable routing, and iterative fine-tuning with catastrophic forgetting mitigation creating sustainable self-improvement.

\section{Objective}

The primary objective of this research is to develop and evaluate CAEM (Confidence-Aware Episodic Memory with Self-Improvement), a system that helps AI models to make fewer hallucinations over time through remembering verified knowledge and iterative refinement. This comprehensive goal decomposes into seven specific objectives:

\begin{itemize}
    \item \textbf{Design a Verified Episodic Memory System with Multi-Layer Quality Control:} Develop an episodic memory architecture that stores question-reasoning-answer tuples only after passing comprehensive verification. Implement a three-layer verification cascade combining natural language inference for factual correctness which checks if the answer matches known facts, self-consistency for logical coherence which checks if the reasoning makes sense and does not contradict itself , and semantic entropy for confabulation detection which checks whether the model is guessing or making things up. Design strategic pruning mechanisms to maintain memory quality as capacity approaches limits, prioritizing high-value episodes means keeping memories based on how often used, high success rate, and relatively recent. Target verification accuracy of 85-90\% using multiple checking methods together.
    
    \item \textbf{Develop Multi-Signal Confidence Estimation for Adaptive Routing Decisions:} Implement a confidence estimation framework combining four complementary signals like token-level probability using geometric mean to see how likely the words are according to the model itself, MC Dropout uncertainty through stochastic forward passes to see how much the answers vary if the model answers the same question multiple times, self-consistency via reasoning chain agreement to check whether different reasoning paths lead to the same answer, and semantic entropy for uncertainty quantification to measure how uncertain or unstable the meaning of the answer is. Calibrate confidence scores using temperature scaling to match empirical accuracy. Enable reliable discrimination between high, medium, and low confidence states to support intelligent routing between retrieval, generation, and external verification.
    
    \item \textbf{Implement Three-Tier Adaptive Routing Based on Similarity and Confidence:} Design a hierarchical routing system that adapts processing strategy depending on whether the question is new than before or similar. Tier 1, if the question is identical to a previous one (similarity greater than 0.95), the system just retrieves the answer from the memory. Tier 2, if the question is somewhat similar (similarity 0.75-0.95), the model generates a new answer using generalization. Tier 3, if the question is new (similarity below 0.75), the system uses full retrieval, generation and verification to ensure accuracy. Optimize routing thresholds to balance accuracy, latency, and computational cost.
    
    \item \textbf{Integrate Cutting-Edge Verification Techniques from Recent Literature:} Implement natural language inference using RoBERTa-Large-MNLI for premise-hypothesis entailment checking on factual claims. The system runs various reasoning attempts and checks to ensure consistency. Incorporate semantic entropy computation following Farquhar et al. (Nature 2024) \cite{farquhar2024semantic}, clustering diverse samples by bidirectional entailment to detect confabulation. Develop dataset-adaptive verification strategies that select primary and secondary verification layers based on task characteristics.
    
    \item \textbf{Create a Closed Self-Improvement Loop over Multiple Learning Cycles:} Establish an iterative refinement process where verification creates high-quality training data, fine-tuning improves model generation accuracy, and improved generation produces even better training data in subsequent cycles. Implement catastrophic forgetting mitigation through mixed training data (90\% verified examples, 10\% general tasks) and L2 regularization. Target progressive accuracy improvement across three cycles while maintaining at least 95\% retention of general language capabilities measured on MMLU benchmark.
    
    \item \textbf{Demonstrate 20-30\% Hallucination Reduction on Established Benchmarks:} Evaluate system performance on four well-known benchmarks like HotpotQA for multi-hop reasoning with explicit evidence, StrategyQA for implicit reasoning requiring world knowledge, FEVER for direct fact verification, and TruthfulQA for resistance to common misconceptions. Track how accurate the answers are and how often hallucinations occur across three learning cycles. Compare against strong baselines including zero-shot prompting, chain-of-thought, retrieval-augmented generation, self-consistency, and vanilla fine-tuning without verification.
    
    \item \textbf{Provide Theoretical Analysis of Convergence and Data Purity Dynamics:} Derive and validate the data purity formula proving that the verified data purity can exceed verification accuracy when base model performance surpasses random chance. Conduct sensitivity analysis across varying verification accuracy levels (75-90\%) to demonstrate robustness. Predict that the system will stabilize after 7-10 learning cycles, achieving 85-90\% high quality memory coverage. And, by measuring actual verified answer quality during experiments, confirm the predictions.
\end{itemize}

These objectives collectively address the three fundamental limitations identified in the Problem Statement. The first three objectives provide episodic memory, confidence estimation, and intelligent routing to overcome stateless operation and poor calibration. Objectives four through six create the self-improvement mechanism with empirical validation. Objective seven provides theoretical foundations ensuring the approach is grounded in principled analysis of verification and learning dynamics.

\section{Methodology in Brief}

This section provides a high-level overview of the CAEM system architecture. In the methodology chapter, detailed technical specifications will be presented.

\subsection{System Architecture Overview}

CAEM integrates five core components. One, the \textit{verified episodic memory} stores 20,000 concrete question-reasoning-answer tuples validated through complete verification, indexed with FAISS for efficient retrieval. Two, the \textit{fine-tuned model} serves as implicit semantic memory, learning generalized patterns from verified examples through iterative training. Three, \textit{Multi-signal confidence estimation} combines token probability, MC Dropout uncertainty, self-consistency, and semantic entropy to produce calibrated confidence scores. Four, \textit{Three-tier adaptive routing} directs queries to retrieval (exact matches), generation (similar patterns), or RAG (novel problems) based on similarity and confidence. Five, \textit{Multi-layer verification} validates all outputs through NLI entailment, self-consistency checking, and semantic entropy analysis before storage.

Figure~\ref{fig:architecture} illustrates the system architecture.

\begin{figure}[p]
\centering
\vspace*{-2cm}
\resizebox{!}{0.95\textheight}{%
\rotatebox{90}{%
\begin{tikzpicture}[
    box/.style={rectangle, draw=#1, line width=2pt, fill=#1, 
                text width=2.3cm, align=center, minimum height=1.6cm, 
                rounded corners=4pt, font=\small\bfseries},
    smallbox/.style={rectangle, draw=#1, line width=1.5pt, fill=#1, 
                text width=1.8cm, align=center, minimum height=1.2cm, 
                rounded corners=3pt, font=\scriptsize\bfseries},
    storage/.style={cylinder, draw=green!60!black, line width=1.8pt, 
                   shape border rotate=90, aspect=0.3, 
                   minimum height=1.5cm, minimum width=2.3cm, 
                   fill=green!8, font=\small\bfseries, align=center},
    external/.style={rectangle, draw=blue!60, line width=1.3pt, fill=blue!5,
                    text width=1.6cm, align=center, minimum height=1cm,
                    rounded corners=2pt, font=\tiny, dashed},
    arrow/.style={->, >=stealth, line width=2pt, #1},
    darrow/.style={->, >=stealth, line width=1.8pt, dashed, #1},
    thindarrow/.style={->, >=stealth, line width=1.3pt, dashed, #1}
]

% Row 1: Input -> Encode -> Route
\node[box=orange!60] (input) at (0,0) {1 INPUT\\[0.15cm]\footnotesize Question};
\node[box=teal!60, right=1.8cm of input] (encode) {2 ENCODE\\[0.15cm]\footnotesize SBERT};
\node[storage, above=0.5cm of encode] (mem) {Memory\\20K};
\node[box=teal!60, right=1.8cm of encode] (route) {3 ROUTE\\[0.15cm]\footnotesize Similarity};

% Tier 1 - Simple retrieval (top)
\node[box=green!60, right=2.5cm of route, yshift=2.5cm] (t1) {
    \textbf{TIER 1}\\
    \scriptsize Retrieve\\
    \scriptsize $>$0.95\\
    \scriptsize 80ms\\
    \scriptsize\textbf{50\%}
};

% Tier 2 - Generation with confidence gate (middle)
\node[box=orange!60, right=1.8cm of route] (t2gen) {
    \textbf{TIER 2}\\
    \scriptsize Generate\\
    \scriptsize 0.75--0.95
};
\node[smallbox=orange!60, right=1.5cm of t2gen] (t2conf) {
    4 CONF\\[0.1cm]
    \tiny Check\\
    \tiny $u>0.6$?
};

% Tier 3 - RAG (bottom)
\node[box=red!60, right=1.8cm of route, yshift=-2.5cm] (t3) {
    \textbf{TIER 3}\\
    \scriptsize RAG\\
    \scriptsize $<$0.75\\
    \scriptsize 3.2s\\
    \scriptsize\textbf{15\%}
};

% External KB for Tier 3
\node[external, left=5.5cm of t3] (extkb) {External\\KB\\(Wiki)};

% Verify and Output
\node[box=red!60, right=2.5cm of t2conf] (verify) {5 VERIFY\\[0.15cm]\footnotesize 85--90\%};
\node[box=green!60, right=1.8cm of verify] (output) {6 OUTPUT\\[0.15cm]\footnotesize Answer};

% Bottom: Model
\node[storage, below=2.5cm of route] (model) {Fine-Tuned\\Model};

% Main flow: Input -> Encode -> Route
\draw[arrow=teal!70] (input) -- (encode);
\draw[darrow=green!70] (encode) -- node[above, font=\scriptsize, fill=white, inner sep=1pt, text=black] {Search} (mem);
\draw[darrow=green!70] (mem) -- node[above, font=\scriptsize, xshift=5mm, yshift=1mm, text=black] {sim score} (encode);
\draw[arrow=teal!70] (encode) -- (route);

% Routing based on similarity
\draw[arrow=green!70] (route) -- node[right, font=\scriptsize, xshift=0.5mm, yshift=1mm, fill=white, inner sep=1pt, text=black] {$s{>}0.95$} (t1);
\draw[arrow=orange!70] (route) -- node[right, font=\scriptsize, yshift=2mm, text=black] {$0.75{<}s{\leq}0.95$} (t2gen);
\draw[arrow=red!70] (route) -- node[right, font=\scriptsize, xshift=0.5mm, yshift=-1mm, fill=white, inner sep=1pt, text=black] {$s{<}0.75$} (t3);

% CRITICAL FIX 1: Memory retrieves to Tier 1
\draw[arrow=green!70, line width=2.5pt] (mem) to[out=-45, in=90] node[right, font=\scriptsize, xshift=2mm, yshift=1mm, fill=white, inner sep=1pt, text=black] {Retrieve} (t1);

% Tier 1: Serves answer directly from memory -- skips Stage 5 entirely
% Freshness is maintained by retroactive re-verification each cycle, not per-query

% CRITICAL FIX 2: Model generates for Tier 2 (SOLID, PROMINENT)
\draw[arrow=orange!70, line width=2.5pt] (model) -- node[left, font=\scriptsize, xshift=-2mm, fill=white, inner sep=1pt, text=black] {Generate} (t2gen);

% Tier 2: Generate -> Confidence Check -> Decision
\draw[arrow=orange!70] (t2gen) -- (t2conf);
\draw[arrow=orange!70] (t2conf) -- node[above, font=\scriptsize, fill=white, inner sep=1pt, text=black] {Accept} (verify);

% CRITICAL FIX 3: Confidence fallback to Tier 3 (MORE PROMINENT)
\draw[arrow=red!70, line width=2pt] (t2conf) -- node[right, font=\scriptsize, xshift=1mm, yshift=-0.5mm, fill=white, inner sep=1pt, text=black] {Fallback} node[right, font=\scriptsize, xshift=1mm, yshift=-3.5mm, fill=white, inner sep=1pt, text=black] {$u{<}0.6$} (t3);

% External KB feeds Tier 3
\draw[darrow=blue!60] (extkb) -- (t3);

% Tier 3: To verify
\draw[arrow=red!70] (t3) -| (verify);

% Output
\draw[arrow=green!70] (verify) -- (output);

% Self-improvement loop
\draw[darrow=purple!70, line width=2.5pt] (verify) |- node[pos=0.7, right, font=\scriptsize, xshift=3mm, text=black] {Store} (mem);
\draw[darrow=purple!70, line width=2.5pt] (mem) |- node[near start, left, font=\scriptsize, xshift=-1mm, yshift=3mm, fill=white, inner sep=1pt, text=black] {Training Data} (model);
\draw[darrow=purple!70, line width=1.8pt] (verify) |- node[near end, right, font=\scriptsize, xshift=1mm, yshift=-2mm, fill=white, inner sep=1pt, text=black] {Train} (model);

% Loop box
\node[draw=purple!70, line width=2.5pt, dashed, rectangle, rounded corners=6pt,
      fit={(model) (mem)}, inner sep=0.5cm,
      label={[font=\normalsize\bfseries, text=purple!70]above:7 SELF-IMPROVEMENT}] {};

% Add note for Tier 2
\node[below=0.1cm of t2gen, font=\scriptsize, text=orange!70!black] {\textit{35\% / 850ms}};

\end{tikzpicture}
}%
}

\vspace{0.3cm}

{\small\caption[Proposed CAEM Architecture]{Proposed CAEM Architecture: 7-stage pipeline with 3-tier routing. Tier 1: direct retrieval from memory (estimated 200-400ms). Tier 2: generate with fine-tuned model, then check confidence to accept or fallback (estimated 1.5-2.5s when successful). Tier 3: RAG for novel/low-confidence queries (estimated 3-6s). Multi-layer verification (target: up to 85--90\%) enables self-improvement through verified knowledge accumulation. \textit{Latencies/confidence values are engineering estimates to be validated empirically.}}}
\label{fig:architecture}
\end{figure}

\subsection{Query Processing and Routing}

When a query comes in, the system turns it into a numerical representation using Sentence-BERT, searches its memory quickly using FAISS and checks how confident it is using four different methods. Routing decisions employ a hierarchical strategy. Tier 1, if the question is almost identical to something already stored (similarity greater than 0.95), just retrieve it from memory. Tier 2, if the question is somewhat similar (similarity between 0.75 and 0.95), the system generates a new answer using fine-tuned model for similar patterns. And, Tier 3, if the question is new (similarity below 0.75), the system uses RAG and full verification for accuracy surety. This progressive strategy achieves efficiency gains as memory coverage expands across cycles.

\subsection{Multi-Layer Verification}

All outputs undergo three-layer verification achieving 85-90\% combined accuracy. Layer 1 uses NLI (Natural Language Inference) to verify factual correctness against retrieved evidence. Layer 2 generates multiple reasoning chains and measures agreement through checking if they all agree. Layer 3 applies semantic entropy from Farquhar et al. (2024) \cite{farquhar2024semantic}, clustering diverse samples by meaning to detect confabulation. Depending on the type of dataset, the system chooses which verification layers are most important to focus on.

\subsection{Self-Improvement Loop}

CAEM implements iterative refinement over three cycles. Each cycle processes questions through the current system, collecting verified episodes that serve as training data. Fine-tuning employs mixed training data (90\% verified, 10\% general tasks) and L2 regularization to prevent catastrophic forgetting. As the model improves, verification becomes more accurate, verified answers create better training data and this leads to steady improvement, targeting a 20--30\% hallucination reduction across three cycles (estimated progression from approximately 60\% to 87\% accuracy on the evaluation benchmarks).

\subsection{Implementation}

The system employs Flan-T5 Large (780M parameters) for its instruction-tuned capabilities and encoder-decoder architecture. Sentence-BERT converts questions into numerical representations, FAISS searches memory quickly for similar questions, and RoBERTa-Large-MNLI performs NLI verification. Evaluation uses HotpotQA, StrategyQA, FEVER, and TruthfulQA benchmarks, measuring exact match accuracy, F1 score, and hallucination rate. The 12-month timeline allocates four months for implementation, four for experimental cycles, and four for analysis and writing.

Figure~\ref{fig:experimental-workflow} illustrates the planned experimental workflow across iterative self-improvement cycles. Each cycle builds upon verified knowledge from previous cycles to drive progressive improvements in accuracy, data purity, and computational efficiency.

\begin{figure}[htbp]
\centering
\begin{tikzpicture}[scale=0.75, transform shape,
    node distance=1.5cm,
    cycle/.style={rectangle, draw, line width=1.3pt, text width=4.5cm, align=left, 
                 minimum height=3cm, rounded corners=4pt, font=\small,
                 drop shadow={shadow xshift=0.3ex, shadow yshift=-0.3ex, opacity=0.4}},
    arrow/.style={->, >=stealth, line width=2pt},
    consolidate/.style={rectangle, draw, line width=1pt, fill=gray!10, 
                       text width=4.5cm, align=center, minimum height=0.8cm, 
                       rounded corners=2pt, font=\footnotesize\itshape}
]

% Cycle 1 - Red to Yellow gradient
\node[cycle, fill=red!20, draw=red!70!black] (c1) {
\textbf{Cycle 1} \textit{(Baseline)}\\
\vspace{0.1cm}
$\bullet$ Memory: \textbf{Empty}\\
$\bullet$ Process: questions\\
$\bullet$ Generate: \textbf{100\%}\\
$\bullet$ Collect verified data\\
$\bullet$ Accuracy: baseline $\rightarrow$ \textbf{improved}
};

% Consolidation 1
\node[consolidate, below=0.3cm of c1] (cons1) {
Fine-tune on verified data\\
Accumulated episodes $\rightarrow$ Model improves
};

% Cycle 2 - Yellow gradient
\node[cycle, fill=yellow!25, draw=orange!70!black, below=0.8cm of cons1] (c2) {
\textbf{Cycle 2} \textit{(Improved)}\\
\vspace{0.1cm}
$\bullet$ Memory: \textbf{Growing}\\
$\bullet$ Process: questions\\
$\bullet$ Retrieve increases\\
$\bullet$ Collect more verified data\\
$\bullet$ Accuracy: \textbf{further improvement}
};

% Consolidation 2
\node[consolidate, below=0.3cm of c2] (cons2) {
Fine-tune on cumulative data\\
More episodes $\rightarrow$ Further improvement
};

% Cycle 3 - Green gradient
\node[cycle, fill=green!25, draw=green!70!black, below=0.8cm of cons2] (c3) {
\textbf{Cycle 3+} \textit{(Target)}\\
\vspace{0.1cm}
$\bullet$ Memory: \textbf{Mature}\\
$\bullet$ Process: questions\\
$\bullet$ High retrieval rate\\
$\bullet$ Sustained quality\\
$\bullet$ Accuracy: \textbf{target performance}
};

% Metrics boxes with better styling
\node[right=1.2cm of c1, text width=3.8cm, align=left, font=\scriptsize, 
      fill=red!10, inner sep=4pt, rounded corners=3pt, draw=red!50] (m1) {
\textbf{Expected (Cycle 1):}\\
$\bullet$ Data purity: 85-95\%\\
$\bullet$ Memory coverage: Low\\
$\bullet$ Avg latency: 2-4s
};

\node[right=1.2cm of c2, text width=3.8cm, align=left, font=\scriptsize, 
      fill=yellow!10, inner sep=4pt, rounded corners=3pt, draw=orange!50] (m2) {
\textbf{Expected (Cycle 2):}\\
$\bullet$ Data purity: improving\\
$\bullet$ Memory coverage: Medium\\
$\bullet$ Avg latency: 1.5-3s
};

\node[right=1.2cm of c3, text width=3.8cm, align=left, font=\scriptsize, 
      fill=green!10, inner sep=4pt, rounded corners=3pt, draw=green!50] (m3) {
\textbf{Target (Cycle 3+):}\\
$\bullet$ Data purity: 88-96\%\\
$\bullet$ Memory coverage: High\\
$\bullet$ Avg latency: 1-2s\\
\vspace{0.1cm}
\textit{\textcolor{green!60!black}{Progressive improvement}}
};

% Arrows with gradient colors
\draw[arrow, red!60!black] (c1) -- (cons1);
\draw[arrow, orange!60!black] (cons1) -- (c2);
\draw[arrow, orange!60!black] (c2) -- (cons2);
\draw[arrow, green!60!black] (cons2) -- (c3);

\end{tikzpicture}
\caption[Experimental Workflow]{Three self-improvement cycles: process 5K questions, collect verified data, fine-tune. Progressive improvement estimated: +15-20\% accuracy.}
\label{fig:experimental-workflow}
\end{figure}

\section{Scopes and Challenges}

\subsection{Research Scope}

This research focuses on tasks where there is a clear right or wrong answer and not subjective questions. The system's target is to handle questions that need to find information from sources, connecting multiple pieces of information, checking facts and avoiding common false beliefs. The four evaluation benchmarks (HotpotQA, StrategyQA, FEVER, TruthfulQA) are chosen because they have questions with definite answers which can be checked objectively. 

The scope explicitly \textit{excludes} creative generation tasks such as story writing, poetry composition, or open-ended dialogue where evaluation criteria are subjective and multiple valid responses exist. Similarly, tasks requiring up-to-date real-world information beyond the training cutoff such as stock prices, breaking news and current events are outside scope, as the system lacks mechanisms for continuous knowledge based updates. Important grounds like Code generation and mathematical theorem proving are excluded to maintain focus on natural language reasoning. The system operates on English text only, with the scope of working on multilingual capabilities in the future.

The model selection reflects resource constraints appropriate for undergraduate thesis work. Flan-T5 Large (780 million parameters) provides sufficient capacity for complex reasoning while remaining feasible for fine-tuning with university GPU allocations. Larger models such as LLaMA-13B or GPT-3.5 scale would require computational resources that are too expensive and resource-heavy. The three-cycle experimental design captures the essential self-improvement dynamics while remaining achievable within a 12-month timeline, though in theory, the convergence requires 7-10 cycles for practical equilibrium.

Instead of using every question in the dataset, the system will be tested on about 5,000 questions per benchmark. This number is enough to get meaningful results but small enough to avoid excessive computational costs. The system can store 20,000 past question-answer-reasoning episodes in memory, which is large enough to test memory management strategies but small enough to run on university-level hardware resources.

\subsection{Technical Challenges}

The implementation of CAEM faces six primary technical challenges that must be addressed:

\begin{itemize}
    \item \textbf{Verification Accuracy Limitations:} The multi-layer verification system targets 85-90\% accuracy, meaning 10-15\% of stored episodes may contain errors. The system tends to produce cleaner data, but any single wrong memory could teach the model a misinformation. The challenge lies in maximizing verification accuracy while maintaining computational feasibility. Semantic entropy verification requires generating ten samples and performing bidirectional NLI classification for clustering, creating substantial overhead. Trade-offs between verification thoroughness and system efficiency must be carefully balanced.
    
    \item \textbf{Catastrophic Forgetting During Fine-Tuning:} Iterative fine-tuning on domain-specific verified examples risks degrading general language capabilities, a phenomenon called \textit{catastrophic forgetting}. While mitigation strategies including mixed training data (90\% verified, 10\% general) and L2 regularization aim to preserve at least 95\% of baseline performance on general tasks measured by MMLU, the effectiveness of these techniques must be validated empirically. If the model still forgets too much, mixing more general data (20\%), efficient fine-tuning methods like LoRA or stopping fine-tuning and just using episodic memory for answers can be some of the possible fixes.
    
    \item \textbf{Confidence Calibration Complexity:} Combining four confidence signals (token probability, MC Dropout, self-consistency, semantic entropy) requires determining optimal weights. At first, we start by treating all four signals equally (0.25 each), but later we can adjust weights based on performance data. However, weight optimization requires sufficient labeled data distinguishing correct from incorrect model outputs, creating a bootstrapping challenge. Another adjustment called Temperature scaling for calibration adds another hyperparameter requiring validation set tuning. Calculating confidence can be expensive, because MC Dropout and semantic entropy need multiple computations per question. So, the system must balance accuracy with speed and cost.
    
    \item \textbf{Memory Retrieval Quality and Pruning Strategy:} The system uses FAISS to quickly find similar past questions, which as introduces potential errors where similar questions are not retrieved despite existing in memory. When the memory is nearly full, the system has to decide which episodes to remove. It considers retrieval frequency (how often episodes are used), success rate (how often retrieved episodes help), confidence (quality at storage time), and recency (newer episodes may be more relevant). Suboptimal pruning could remove valuable episodes while retaining redundant or low-quality entries. The novelty detection threshold of 0.95 similarity requires careful tuning, for example, too high allows many duplicates and too low rejects genuinely novel variations.
    
    \item \textbf{Baseline Performance Dependencies:} The theoretical proof assumes that the baseline model has learned meaningful patterns and performs better than random guessing, meaning accuracy above 50\% for binary classification and above 25\% for four-way. If Flan-T5 Large achieves below 50\% accuracy on specific benchmark subsets, storing verified answers could actually make things worse. To reduce this risk, the project chooses Flan-T5 Large, which can follow instructions well and picks benchmarks where the model already does reasonably well.
    
    \item \textbf{Evaluation Metric Selection:} Measuring hallucination reduction requires careful metric design beyond simple accuracy. The FActScore methodology for atomic fact verification provides a gold standard but it takes a lot of effort, requiring manual annotation or expensive API calls to verification models. So, there is a balance between checking everything perfectly and keeping evaluation doable. Binary accuracy metrics may miss subtle improvements in explanation quality or hedge expressions. Comparing against strong baselines requires faithful reimplementation, as published results often differ in how they were measured or which data they used.
\end{itemize}

\subsection{Mitigation Strategies}

For each challenge identified, the research design incorporates mitigation strategies. The three layer verification system, which is NLI, self-consistency , semantic entropy, ensures answers are checked carefully which will reduce errors. Catastrophic forgetting monitoring includes early stopping criteria. For example, if MMLU score drops below 93\% of baseline, training halts and the previous cycle's model is retained. Confidence calibration begins with equal weighting, deferring optimization to future work if validation data proves insufficient. The system removes old and less-valuable memories only when reaching 95\% capacity and value-based metrics combining multiple criteria like usage, success rate and recency. Computational constraints are managed through efficient implementation choices (DeepSpeed, gradient checkpointing) and scoped experimentation. Before full experiments, the system checks all baseline methods to make sure if everything works.

Although the problems are significant, they can be managed within the resources and the time frame. The aim is to improve the system clearly and measurably while focusing on using smart design choices to reduce errors.